\subsection{Divisors attached to torsion modules}

We keep the same notation and hypotheses as in nos. 2, 3 and 4. Recall that D(A) (or simply D) denotes the divisor group of A, written additively: we know (§ 1, no. 3, Theorem 2) that D is the free Z-module generated by the elements of P.

Let T be a finitely generated torsion A-module. For all \(p \in P\) , \(T_{p}\) is a finitely generated torsion \(A_{p}\) -module and hence a module of finite length (Chapter IV, §2, no. 5, Corollary 2 to Proposition 7); we shall denote this length by \(l_{p}(T)\) . Now \(T_{p} = 0\) for all p not containing the annihilator of T and hence for almost all p (§1, no. 6, Theorem 4), which justifies the following definition:

DEFINITION 4. If T is a finitely generated torsion A-module, the divisor:

\[
\chi (\mathbf {T}) = \sum_ {\mathfrak {p} \in \mathbf {P}} l _ {\mathfrak {p}} (\mathbf {T}). \mathfrak {p}.
\]

is called the content of T.

PROPOSITION 10. (i) Let \(0 \to T_1 \to T_2 \to T_3 \to 0\) be an exact sequence of finitely generated torsion A-modules. Then

\[
\chi (\mathrm{T} _ {2}) = \chi (\mathrm{T} _ {1}) + \chi (\mathrm{T} _ {3}).
\]

\begin{enumerate}
\def\labelenumi{(\roman{enumi})}
\setcounter{enumi}{1}
\item
  If there exists a pseudo-isomorphism \(f \colon \mathrm{T}_1 \to \mathrm{T}_2\) , then \(\chi(\mathrm{T}_1) = \chi(\mathrm{T}_2)\) .
\item
  In order that \(\chi(T) = 0\) , it is necessary and sufficient that \(T\) be pseudo-zero.
\end{enumerate}

In view of Definition 4, it suffices to consider for each \(p \in P\) the values of \(l_{p}\) for the torsion modules considered. Property (i) then follows from Chapter II, §2, no. 4, Theorem 1 and the additivity of lengths in an exact sequence (Algebra Chapter II, §1, no. 10, Proposition 16) and properties (ii) and (iii) follow immediately from the definitions in no. 4.

COROLLARY. Let \(0 \to T_n \to T_{n-1} \to \cdots \to T_0 \to 0\) be an exact sequence of finitely generated torsion A-modules. Then \(\sum_{i=0}^{n} (-1)^i \chi(T_i) = 0\) .

In view of Chapter II, §2, no. 4, Theorem 1, this follows again from the analogous property of the \(l_{\mathfrak{p}}\) (Algebra, Chapter II, §1, no. 10, Corollary 3 to Proposition 16).

Recall (Chapter II, § 5, no. 4) that we may speak of the set \(\mathbf{F}(\mathbf{A})\) of classes of finitely generated A-modules with respect to the relation of isomorphism; for every finitely generated A-module M, let \(\operatorname{cl}(\mathbf{M})\) denote the corresponding element of \(\mathbf{F}(\mathbf{A})\) ; we shall denote by \(\mathbf{T}(\mathbf{A})\) the subset of \(\mathbf{F}(\mathbf{A})\) consisting of the classes of finitely generated torsion A-modules. Clearly \(\chi\) defines a mapping of \(\mathbf{T}(\mathbf{A})\) to \(\mathbf{D}(\mathbf{A})\) , also denoted by \(\chi\) , such that \(\chi(\operatorname{cl}(\mathbf{T})) = \chi(\mathbf{T})\) .

PROPOSITION 11. Let G be a commutative group written additively and \(\phi: T(A) \to G\) a mapping; for every finitely generated torsion A-module T, we also write, by an abuse of language, \(\phi(T) = \phi(\text{cl}(T))\) . Suppose that the following conditions are satisfied:

\begin{enumerate}
\def\labelenumi{(\arabic{enumi})}
\item
  If \(0 \to \mathrm{T}_1 \to \mathrm{T}_2 \to \mathrm{T}_3 \to 0\) is an exact sequence of finitely generated torsion A-modules, then \(\phi(\mathrm{T}_2) = \phi(\mathrm{T}_1) + \phi(\mathrm{T}_3)\) .
\item
  If \(\mathbf{T}\) is pseudo-zero, then \(\phi (\mathbf{T}) = 0\) .
\end{enumerate}

Then there exists a unique homomorphism \(\theta: \mathbf{D}(\mathbf{A}) \to \mathbf{G}\) such that \(\phi = \theta \circ \chi\) .

As \(\chi(A/p) = p\) for all \(p\) , necessarily \(\theta(p) = \phi(A/p)\) for all \(p \in P\) , which proves the uniqueness of \(\theta\) , since the elements of \(P\) form a basis of \(D(A)\) . Conversely, let \(\theta\) be the homomorphism from \(D(A)\) to \(G\) such that \(\theta(p) = \phi(A/p)\) for all \(p \in P\) and let us show that it solves the problem. For this, we write

\[
\psi (\mathbf {T}) = \phi (\mathbf {T}) - \theta (\chi (\mathbf {T}))
\]

for every finitely generated torsion A-module T; clearly conditions (1) and (2) are also satisfied if \(\phi\) is replaced by \(\psi\) . On the other hand, \(\psi(A/p) = 0\) if \(p \in P\) ; if \(p\) is a prime ideal \(\neq 0\) and not in \(P\) , the annihilator of \(A/p\) is contained in no ideal of \(P\) , hence (no. 4, Theorem 5) \(A/p\) is pseudo-zero and therefore \(\psi(A/p) = 0\) . This being so, every finitely generated torsion A-module T admits a decomposition series whose factors are isomorphic to A-modules of the form \(A/p\) , where \(p \in \text{Supp}(T)\) (Chapter IV, § 1, no. 4, Theorems 1 and 2), and hence \(p \neq 0\) since \(T\) is a torsion module. By induction on the length of this decomposition series, we deduce (in view of property (1) for \(\psi\) ) that \(\psi(T) = 0\) .

* We may, as in no. 4, consider the quotient category \(\mathcal{T} / \mathcal{T}'\) of the category \(\mathcal{T}\) of finitely generated torsion A-modules by the full sub-category \(\mathcal{T}'\) of pseudo-zero finitely generated torsion A-modules. In the language of Abelian categories, Proposition 11 then expresses the fact that the Grothendieck group of the Abelian category \(\mathcal{T} / \mathcal{T}'\) is canonically isomorphic to D(A). *

PROPOSITION 12. If a is an ideal \(\neq 0\) of A,

\[
\chi (\mathrm{A} / \mathfrak {a}) = \chi ((\mathrm{A}: \mathfrak {a}) / \mathrm{A}) = \operatorname{div} \mathfrak {a}.
\]

Let \(p \in P\) . Then \(aA_p = p^{n_p}A_p\) where \(n_p \geqslant 0\) , since \(A_p\) is a discrete valuation ring. As \((A/a)_p = A_p/aA_p\) , \(l_p(A/a) = n_p\) , whence \(\chi(A/a) = \sum_{p \in P} n_p p = \text{div } a\) ( \(\S 1, \text{ no. } 4, \text{ Proposition } 7\) ).

On the other hand, \((\mathbf{A}:\mathfrak{a})_{\mathfrak{p}} = \mathbf{A}_{\mathfrak{p}}\colon \mathfrak{a}\mathbf{A}_{\mathfrak{p}} = \mathfrak{p}^{-n_{\mathfrak{p}}}\mathbf{A}_{\mathfrak{p}}\) , hence \(l_{\mathfrak{p}}((\mathbf{A}:\mathfrak{a}) / \mathbf{A}) = n_{\mathfrak{p}}\) and we conclude in the same way.

